% !TeX program = lualatex
\documentclass[11pt]{article}
\usepackage[margin=25mm]{geometry}
\usepackage{fontspec}
\setmainfont{Latin Modern Roman}
\usepackage{amsmath,amssymb,mathtools}
\usepackage{unicode-math}
\setmathfont{Latin Modern Math}
\usepackage{xurl,hyperref}
\hypersetup{colorlinks=true,urlcolor=blue}
\setcounter{secnumdepth}{0}
\setlength{\parindent}{0pt}
\setlength{\parskip}{5pt}
\setlength{\emergencystretch}{3em}
\newcommand{\R}{\mathbb R}
\newcommand{\N}{\mathbb N}
\newcommand{\Z}{\mathbb Z}
\newcommand{\Q}{\mathbb Q}
\newcommand{\C}{\mathbb C}
\newcommand{\D}{\mathbb D}
\newcommand{\F}{\mathbb F}
\newcommand{\sref}[2]{\hyperlink{source:#1:#2}{[S#2]}}
\newcommand{\cataloguescope}[1]{#1}
\begin{document}
\section*{Invariant anticanonical sections under structure-sheaf acyclicity}\label{30007064}
\subsection{Definitions and mathematical statement}
\paragraph{Statement recorded in the supplied source.}
For a smooth connected projective complex X with smoothly semipositive $-K_X$ and $H^i(X,O_X)=0$ for all $i>0$, must every algebraic torus action admit a nonzero invariant section of some positive anticanonical power with its natural linearization?
\subsection{Short English statement}
For a smooth connected projective complex X with smoothly semipositive $-K_X$ and $H^i(X,O_X)=0$ for all $i>0$, must every algebraic torus action admit a nonzero invariant section of some positive anticanonical power with its natural linearization?
\subsection{Sources}
\begin{itemize}
\item[S1] OpenAI, manuscript citation page. Source of the formulation motivating this local problem.
\url{https://github.com/openai/math/blob/adc7f1241b42e322a6451854ab7e4b4c146bf78a/preprints/Exact-orders-and-invariant-anticanonical-linear-systems-September-26-2026/README.md}.
\end{itemize}
\end{document}
